\documentclass[10pt,a4paper]{amsart}
\usepackage[margin=19mm]{geometry}
\usepackage{amsmath,amssymb,mathtools}
\usepackage[scr=rsfs,cal=euler]{mathalfa}
\usepackage{xfrac}
\usepackage[unicode,hidelinks,bookmarks=false]{hyperref}
\newcommand{\ie}{i.e.\ }
\newcommand{\cf}{cf.\ }
\newcommand{\resp}{respectively\ }
\newcommand{\Subsection}{\subsection}
\providecommand{\nobreakdash}{\nobreak\hbox{-}}
\setlength{\emergencystretch}{2em}
\multlinegap0pt
\usepackage{bm}
\usepackage{bbm}
\usepackage{dsfont}
\usepackage{xspace}
\usepackage{cancel}
\usepackage{mathtools}
\usepackage{amsthm}
\makeatletter
\@ifundefined{theorem}{\newtheorem{theorem}{Theorem}}{}
\@ifundefined{example}{\newtheorem{example}[theorem]{Example}}{}
\@ifundefined{lemma}{\newtheorem{lemma}[theorem]{Lemma}}{}
\@ifundefined{proposition}{\newtheorem{proposition}[theorem]{Proposition}}{}
\makeatother
\makeatletter
\expandafter\ifx\csname LIST\endcsname\relax
\newenvironment{LIST}[1]{%
\setlength{\tmplenghta}{#1}
\setlength{\tmplenghtb}{#1}
\setlength{\tmplenghtc}{#1}
\advance\tmplenghtb-5pt
\advance\tmplenghtc 42pt
\setcounter{tmpc}{0}
\begin{list}{{\rm (\alph{tmpc})}}{\usecounter{tmpc}
\setlength{\leftmargin}{\tmplenghta}
\setlength{\rightmargin}{0cm}
\setlength{\itemsep}{1pt}
\setlength{\topsep}{3pt}
\setlength{\labelsep}{5pt}
\setlength{\labelwidth}{\tmplenghtb}
\setlength{\listparindent}{\tmplenghta}}
}{\end{list}}
\fi
\makeatother
\title[Projective limits in Euclidean quantum field theory, I: Free]{Projective limits in Euclidean quantum field theory, I: Free scalar fields}
\author{Svetoslav Zahariev}
\date{Source: arXiv:2509.12498v1 (CC BY 4.0)}
\begin{document}
\maketitle

\noindent\emph{Source and licence.} Projective limits in Euclidean quantum field theory, I: Free scalar fields. Version arXiv:2509.12498v1 (CC BY 4.0), distributed under the Creative Commons Attribution 4.0 International licence, as recorded in the arXiv licence metadata for this exact version and shown on the abstract page https://arxiv.org/abs/2509.12498v1. Full rights notice: licenses/arxiv-2509.12498-CC-BY-4.0.txt.\par\smallskip

\noindent\textit{Editorial scope.} Counted unit: the Theorem whose source label is \texttt{mainprojsythe} in the article (source0.tex lines 291--293, source label \texttt{mainprojsythe}), delivered together with the article's own complete original proof of it (source0.tex lines 294--313, one \texttt{proof} environment of 20 lines). No mathematics is written, repaired or re-derived: statement and proof are the article's own text line for line. Required context retained uncounted: funcpropcharfu (lines 177--179); prooncharfu (lines 190--193); gaumesexam (lines 198--216); cochamapcoij (lines 269--271); chmeasulem (lines 279--287). Every definition, display and label of those lines is retained unchanged. Omitted from this package and not redistributed: 1--176, 180--189, 194--197, 217--268, 272--278, 288--290, 314--452. Nothing inside a retained excerpt is omitted or altered: every retained body is delivered exactly as the article's lines are written, so no omission receipt of any kind accompanies this package. Citations: the retained excerpts carry no citation command at all, so no bibliography entry of the article is transcribed and no reference is redistributed. Reference adaptation: none was needed and none was made; every reference inside the delivered unit resolves inside it, so no unresolved reference or citation is produced. Printed slips kept literal: no printed word, symbol, hypothesis or number of the counted statement or of its proof was altered. Layout and class: the article's own class is \texttt{amsart} with packages including amsaddr, graphicx, amsthm, amssymb, amsmath, calc, eucal, ifthen, amscd, xy, xcolor, pgf, tikz, pgfplots; the wrapper is the lane's pinned \texttt{amsart} editorial wrapper at 10pt on A4, which restates the theorem-like environments the retained text needs and adds the article's own macro definitions that the retained text needs (none), with extra wrapper packages (bm, bbm, dsfont, xspace, cancel, mathtools). The delivered numbering is the unit's own. Assets: no retained span contains a figure environment, a graphics-inclusion command, a PDF-page inclusion, a table or a TikZ picture; no image file is copied into this package, the package declares no asset dependency and no image file is redistributed with it. Licence: version arXiv:2509.12498v1 of this article is distributed under the Creative Commons Attribution 4.0 International licence, as recorded in the arXiv licence metadata for this exact version and shown on the abstract page https://arxiv.org/abs/2509.12498v1. A full rights notice is delivered at licenses/arxiv-2509.12498-CC-BY-4.0.txt. Characters: every retained excerpt is pure ASCII, so the delivered engine, XeLaTeX with its default Latin Modern text font, reports no missing character.
\begin{lemma}\label{funcpropcharfu}
Let $\mathcal{H}_1$ and $\mathcal{H}_2$ be two  real finite dimensional Hilbert spaces and $L: \mathcal{H}_1 \rightarrow \mathcal{H}_2$ be a linear map. Let $\mu_i$ be a finite Borel measure on $\mathcal{H}_i$, $i=1,2$. Then $L(\mu_1)=\mu_2$ if and only if $S_{\mu_1}L^*=S_{\mu_2}$.
\end{lemma}

\begin{proposition}\label{prooncharfu} (a) Assume that the functionals $S_iP_i$ defined on $\mathcal{H}_{\infty}$ are pointwise uniformly bounded and equicontinous. Then there exists a unique continuous positive definite functional $S_{\infty}$ on $\mathcal{H}_{\infty}$ such that $S_{\infty}I_i=S_i$ for all $i$.
	
(b) The triple $(\mathcal{H}_i,I_{ij},S_{\mu_i})$ is an inductive system of characteristic functionals if and only if $(\mathcal{H}_i,I_{ij}^{*},\mu_i)$ is a  projective system of measures.
\end{proposition}

\begin{example}\label{gaumesexam}
Let $\mu_{A_i}$ be the centered Gaussian probability measure on $\mathcal{H}_i$ with covariance operator $A_i$ so that one has 
$$ S_{\mu_{A_i}}(h)=e^{-\frac{\langle A_i h, h\rangle_i}{2}}, 
	\quad h \in \mathcal{H}_i.$$
Then one easily sees that $(\mathcal{H}_i,I_{ij},S_{\mu_{A_i}})$ is an inductive system of characteristic functionals if and only if 
\begin{equation}\label{consgaucova}
I_{ij}^{*} A_j I_{ij}=A_i
\end{equation}
holds for all $j>i$.

Now assume in addition that the operators $A_i$ are uniformly bounded, i.e. there exists $K>0$ such that $\|A_i \|_i \leq K $ for all $i$. Then it is not hard to see that the functionals $S_{\mu_{A_i}}P_i$ are equicontinuous. Indeed, a short computation yields
$$ \int_{\mathcal{H}_i}|e^{i\langle h_1, h \rangle_{i}}
-e^{i\langle h_2, h \rangle_{i}} |^2 d\mu_{A_i}(h)=
2(1-e^{-\|A^{1/2}_{i}(h_1-h_2)\|_{i}^{2}/2})$$
for all $h_1,h_2 \in \mathcal{H}_i$. It follows that 
$$| S_{\mu_{A_i}}(P_i(h_1))-S_{\mu_{A_i}}(P_i(h_2))| \leq 
\sqrt{2}(1-e^{-\|A^{1/2}\|_{i}\|h_1-h_2\|^{2}/2})^{1/2}$$
for all $h_1,h_2 \in \mathcal{H}_{\infty}$, which implies the desired equicontinuity. Thus by Proposition \ref{prooncharfu}(a) there exists a unique continuous positive definite functional $S_{\infty}$ on $\mathcal{H}_{\infty}$ such that $S_{\infty}I_i=S_{\mu_{A_i}}$ for all $i$.
\end{example}

\begin{equation}\label{cochamapcoij}
P^1_{ij}d_{i,0} =d_{j,0} P^0_{ij}, \quad i>j.
\end{equation}

\begin{lemma}\label{chmeasulem}
Let $\mathcal{V}_1$ and $\mathcal{V}_2$ be two finite dimensional real vector spaces and let $L$ be a linear isomorphism. Let $f: \mathcal{V}_2 \to \mathbb{R}$ be a positive continuous integrable function.

Define Borel probability measures on $\mathcal{V}_i$
via 
\[ d\mu_1(v_1) = N_1 f(L(v_1)) \, dv_1, \]
\[ d\mu_2(v_2) = N_2 f(v_2) \, dv_2, \]
where $dv_i$ stands for the Lebesgue measure on $\mathcal{V}_i$ and $N_i$ is a normalization constant. Then $\mu_1 = L^{-1}(\mu_2)$.
\end{lemma}

\begin{theorem}\label{mainprojsythe}
The triple $(C^{0,0}(K_i),P^0_{ij},\mu_{i,r})$ is a projective system of measures.
\end{theorem}

\begin{proof}
We write $d_{i,0}^{-1}$ for the inverse of the linear isomorphism from $C^{0,0}(K_i)$ to $\mathtt{Im}\hspace{1pt}d_{i,0}$ given by $d_{i,0}$ and $\nu_i$ for the centered Gaussian probability measure on $\mathtt{Im}\hspace{1pt}d_{i,0}$ with characteristic functional
$$S_{\nu_i} (c)=e^{-\frac{1}{2}\langle c, c \rangle_{1,i}^r},\quad c \in \mathtt{Im}\hspace{1pt}d_{i,0}.$$
Using Lemma \ref{chmeasulem}, we conclude that $\mu_{i,r}=d_{i,0}^{-1}(\nu_i)$, therefore one has
\begin{equation}\label{charconide}
S_{\mu_{i,r}} =S_{\nu_i}(d_{i,0}^{*,r})^{-1}
\end{equation}
by Lemma \ref{funcpropcharfu}. As observed above, the maps $(P^1_{ij})^{*,r}$ are isometries, hence (cf. Example \ref{gaumesexam}) $(\mathtt{Im}\hspace{1pt}d_{i,0},(P^1_{ij})^{*,r}, S_{\nu_i})$ is an inductive system of characteristic functionals, i.e. one has
\begin{equation}\label{twocompcharco}
	S_{\nu_{i}}(P^1_{ij})^{*,r} = S_{\nu_{j}}.
\end{equation}
Further, we infer from (\ref{cochamapcoij}) that
\begin{equation}\label{conjchainma}
(d_{i,0}^{*})^{-1}(P^0_{ij})^{*,r} 
= (P^1_{ij})^{*,r}(d_{j,0}^{*})^{-1}, \quad i>j.
\end{equation}
It follows immediately from (\ref{charconide}), (\ref{twocompcharco}) and (\ref{conjchainma}) that
$ S_{\mu_{i,r}}(P^0_{ij})^{*,r} = S_{\mu_{j,r}}$, which by Lemma \ref{funcpropcharfu} implies that 
$P^0_{ij}(\mu_{i,r})=\mu_{j,r}$.
\end{proof}

\end{document}
